%! Operations with Polynomials — Unit 4, Lesson 4.1 — Warm-Up (Answer Key)
\documentclass[10pt]{article}
\usepackage{algebra2-article}
\usepackage{algebra2-key}   % pulls in -boxes; do NOT also load -boxes
\newcommand{\TallMath}[1]{$\displaystyle #1\rule[-1.4em]{0pt}{3.2em}$}

\begin{document}

\pageheader{Unit 4, Lesson 4.1}{Warm-Up --- Answer Key}
\namedateperiod

\vspace{0.10in}

\begin{notesbox}{Getting Ready: Skills We'll Use Today}
\small These four skills are the whole engine of today's lesson. Work quickly.

\vspace{5pt}
\textbf{1. Multiplying powers.} Multiply the coefficients; do something else with the exponents.

\vspace{2pt}
$x^3 \cdot x^5 = $ \ans{$x^8$} \qquad $(2x^2)(-4x^3) = $ \ans{$-8x^5$} \qquad
$(-3x)(5x^4) = $ \ans{$-15x^5$}

\vspace{3pt}
Complete the rule: to multiply, \textbf{multiply} the coefficients and \ans{add} the exponents.

\vspace{6pt}
\textbf{2. Combining like terms.} Only terms with the \emph{same variable and same exponent} combine.

\vspace{2pt}
$6x^3 - 2x + 5 + x^3 + 7x - 9 = $ \ans{$7x^3 + 5x - 4$}

\vspace{3pt}
Circle the true statement: \quad $3x^2 + 3x^3 = 6x^5$ \qquad $3x^2 + 3x^3$ cannot be combined
\ans{$\leftarrow$ true}

\vspace{6pt}
\textbf{3. Distributing --- watch the negative.}

\vspace{2pt}
$-(4x^2 - 7x + 3) = $ \ans{$-4x^2 + 7x - 3$}

\vspace{3pt}
$3x(2x^2 - 5x + 1) = $ \ans{$6x^3 - 15x^2 + 3x$}

\vspace{6pt}
\textbf{4. A product you already know, and a name from Lesson 4.0.}

\vspace{2pt}
$(x+3)(x-5) = $ \ans{$x^2 - 2x - 15$}

\vspace{3pt}
For $x^3 - 3x + 2$: \; degree $=$ \ans{3} \quad leading coefficient $=$ \ans{1} \quad
so both arms: \ans{left down, right up}

\end{notesbox}

\begin{teachernote}
Item~1 is the engine of every product today (coefficients multiply, exponents \emph{add}) --- the
mirror-image error, adding exponents when \emph{adding} terms, is what item~2 guards against. Item~3
is the single highest-value debrief: ask what happened to the sign of \emph{each} of the three terms,
not just the first, because $-(4x^2-7x+3)$ is exactly the move students will drop in today's
subtraction problems. Item~4's polynomial is Lesson~4.0's anchor $g(x)=x^3-3x+2$; students met it as
$(x-1)^2(x+2)$, and today's Hook asks whether those two forms are really the same function.
\end{teachernote}

\end{document}
